\documentclass{exam}
\usepackage{graphicx} % Required for inserting images
\usepackage{algorithmicx}
\usepackage{algpseudocode}
\usepackage{geometry}[border=1in]
\usepackage{algorithm}
\usepackage{amsmath}
\usepackage{amssymb}
\usepackage{amsthm}
\usepackage{listings}
\usepackage{mathtools}
\usepackage{hyperref}
\DeclarePairedDelimiter\ceil{\lceil}{\rceil}
\DeclarePairedDelimiter\floor{\lfloor}{\rfloor}

\newtheoremstyle{prob}% name 
{7pt}% Space above 
{7pt}% 
{}% 
{0pt}%Indent amount 1 
{\bfseries}% Theorem head font 
{.}% 
{\newline}% 
{}%

\theoremstyle{prob}
\newtheorem*{probstate}{PROBLEM}

\printanswers

\title{Homework 6}
\author{COMP221 Spring 2026 - Suhas Arehalli}
\date{}

\begin{document}

\maketitle

Complete the problem below. Check the course website \& syllabus for further instructions.

If any problem is unclear, or you think you found a typo, please let me know ASAP so I can clarify!

\section*{Problems}

\begin{questions}
    \question \textbf{Degree-Constrained Spanning Trees} \textit{(Based on Skiena 11-12)}
    
    \begin{parts}
        \part The low-degree spanning tree problem (\textproc{Low-Degree-Span}) is a decision problem defined as follows:
        \begin{probstate}[\textproc{Low-Degree-Span}]
            \textbf{Input:} An undirected graph $G = (V, E)$ and $k \in \mathbb{Z}_{\geq 0}$
        
            \textbf{Output:} TRUE iff there exists a spanning tree where each vertex in the tree has degree $\leq k$. FALSE otherwise.
        \end{probstate}
        Prove (using a reduction) that \textproc{Low-Degree-Span} is NP-hard. Make sure you provide both a description of the reduction and a brief justification of the correctness of the reduction as part of your answer.

        \textbf{**Hint**:} Skim through the NP-hard problems you've seen so far (including those in the textbook we didn't discuss in class!). Which involve trees with constrained degree? This might require some clever thinking!
        
        \part Consider the \textit{high-degree spanning tree problem} (\textproc{High-Degree-Span}) which is defined as follows:
        \begin{probstate}[\textproc{High-Degree-Span}]
            \textbf{Input:} An undirected graph $G = (V, E)$ and $k \in \mathbb{Z}_{\geq 0}$.
            
            \textbf{Output:} TRUE if $G$ contains a spanning tree which contains a vertex with degree $\geq k$ (within the spanning tree). FALSE otherwise.
        \end{probstate} 
        Show that \textproc{High-Degree Span} is in P.

        \textbf{**Hint**}: Think through what graph methods we've seen that create spanning trees. Can one force 
        
    \end{parts}
\end{questions}

\end{document}